\documentclass[a4paper,10pt,BCOR12mm]{report}
\NeedsTeXFormat{LaTeX2e}

\usepackage[USenglish]{babel}

\usepackage{times}
\usepackage{graphicx}
\usepackage{amsmath}
\usepackage{amssymb}
\usepackage{listings}

\usepackage{ifpdf}
%\ifpdf
\usepackage[pdfauthor={ZIH, TU Dresden},pdftitle={libSplash User Manual},colorlinks={true},linkcolor={blue},filecolor={blue},pagecolor={blue},urlcolor={blue}]{hyperref}
%\fi

\usepackage{url}

\renewcommand{\chaptermark}[1]{\markboth{\thechapter\ ~#1}{}}
\renewcommand{\sectionmark}[1]{\markright{\thesection\ ~#1}}

\renewcommand{\familydefault}{\sfdefault}

\newcommand{\command}[1]{\small \texttt{#1}}

\begin{document}

\thispagestyle{empty}

\begin{figure*}[ht]
\begin{minipage}[b]{0.45\linewidth}
	\centering
	\noindent\includegraphics[width=5cm]{zih_logo}\\
\end{minipage}
\hspace{0.05\linewidth}
\begin{minipage}[b]{0.45\linewidth}
	\centering
	\noindent\includegraphics[width=5cm]{hzdr_logo}\\
\end{minipage}
\end{figure*}

\vspace{1cm}

{\bfseries\Huge
%\rule[0.75cm]{\textwidth}{1pt}
\noindent
libSplash \\[.5cm]
User Manual\\
\rule{0pt}{0.75cm}\rule{\textwidth}{1pt}
}

{\centering
Last update: \today
}

\vspace{1cm}

{\noindent}TU Dresden\\
Center for Information Services and \\High Performance Computing (ZIH)\\
01062 Dresden\\
Germany\\[1ex]
\url{http://www.tu-dresden.de/zih}\\

\vspace{1cm}

{\noindent}Helmholtz Zentrum Dresden Rossendorf\\
Bautzner Landstrasse 400\\
01328 Dresden\\
Germany\\[1ex]
\url{http://www.hzdr.de}\\

\tableofcontents

\newpage

%-----------------------------------------------------------

\chapter{Introduction}

\section{About libSplash}

libSplash is a combined project of the Center for Information Services and HPC (ZIH) of the Technical University of Dresden and the Helmholtz-Zentrum Dresden-Rossendorf (HZDR).
The project aims at developing a HDF5-based I/O library for HPC simulations. It is created as an easy-to-use frontend for the standard HDF5 library with support for MPI processes in a cluster environment.
While the standard HDF5 library provides detailed low-level control, libSplash simplifies tasks commonly found in large-scale HPC simulations, such as iterative computations and MPI distributed processes.

\section{About this Manual}

This manual describes the general ideas and usage modes of libSplash and its most important classes.
For a detailed explanation of all available classes, interfaces and methods, please refer to the
Doxygen HTML documentation.

\newpage

\section{Installation}

\subsection{Requirements}

Building libSplash from source requires to checkout the trunk via Subversion using the command\\
\command{svn co https://fusionforge.zih.tu-dresden.de/svn/datacollector/trunk/}
Building libSplash requires HDF5 in version 1.8.6 or higher.
To use the CMakeLists.txt file which comes with the source code,
you must have CMake version 2.6 installed.

\subsection{Compiling}

To compile libSplash, it is recommended to create a new \command{build} folder
and execute \command{cmake CMAKE\_INSTALL\_PREFIX=\textless INSTALL\_PATH \textgreater \textless CMakeLists.txt\_PATH \textgreater}\\
Than run \command{make install} to build and install libSplash library and binaries to the chosen installation directory.
This will build libSplash as an shared object library (\emph{libsplash.so}) as well as all tools.

\subsection{Linking}

To use libSplash in your project, you must link against the created shared object library
\emph{libsplash.so}.


\section{Usage}

libSplash consists of two parts: the basic \texttt{DataCollector} interface and the extended \texttt{DomainCollector} interface.

\texttt{DataCollector} is the basic interface for most operations, such as accessing files as well as reading and writing Datasets and Attributes.
This interface is implemented in the \texttt{SerialDataCollector} class.

\texttt{DomainCollector} extends this interface with operations on \emph{Domains} which can represent the simulation area,
a logical data field or a similar logical program structure. This is helpful to allow easy post-mortem access to the stored data from
an analysis or visualization tool.

%-----------------------------------------------------------

\chapter{SerialDataCollector}

\section{Files}

libSplash stores data in HDF5 files with the extension \emph{.h5}.
The filename structure is \emph{(common name)\_(mpi position).h5}.
\emph{common name} is the name chosen by the user for the libSplash files, e.g.
'simulation\_data'.
\emph{mpi position} is the three-dimensional position of the MPI process creating this
file, starting at (0, 0, 0).
This format is chosen even if no MPI environment is used to create the files.

Example:
\begin{itemize}
	\item If libSplash is used from a non-parallel program using one process, only the file \emph{simulation\_data\_0\_0\_0.h5} is created.

	\item If libSplash is used by a MPI parallel program with 2x2 processes, the files
	\emph{simulation\_data\_0\_0\_0.h5}, \emph{simulation\_data\_1\_0\_0.h5},
	\emph{simulation\_data\_0\_1\_0.h5} and \emph{simulation\_data\_1\_1\_0.h5} are created.
\end{itemize}

\subsection{File Structure}

Each libSplash file uses a similar internal file structure which is composed from groups (folders),
datasets, attributes and references.

\begin{itemize}
	\item \textbf{header}
	This group stores general information about this file and the creation context (e.g.
	the number of MPI processes participating in creating all related files).

	\item \textbf{data}
	This group stores the actual (simulation) data and their annotated attributes. Indexed
	sub-groups are used to reflect an iterative program pattern. In the following example, every 10th
	iteration is stored using libSplash.
	\begin{itemize}
		\item \textbf{0}
		Iteration 0
		\begin{itemize}
			\item dataset\_A
			\item dataset\_B
		\end{itemize}

		\item \textbf{10}
		Iteration 10
		\begin{itemize}
			\item dataset\_A
			\item dataset\_B
		\end{itemize}

		\item ...
	\end{itemize}

	\item \textbf{common}
	???
\end{itemize}


\subsection{Opening Files}

Before data can be stored or read, files must be opened by calling
\texttt{DataCollector::open}. This method requires the common part of the filename
and an object of type \texttt{FileCreationAttr}.
This object defines the file access type as well as the number of
participating MPI processes, the MPI position of the calling process and further information.
The following file access types are available:
\begin{itemize}
	\item \texttt{FAT\_CREATE}
	A new file is created. Any existing file with this name if overwritten.
	\texttt{FileCreationAttr} is used to determine the MPI position part of the filename.

	\item \texttt{FAT\_WRITE}
	An existing file is opened for reading and writing.
	\texttt{FileCreationAttr} is used to determine the MPI position part of the filename.
	If the file does not exist, it is created.
	Otherwise, any write access to existing datasets will overwrite them.

	\item \texttt{FAT\_READ}
	An existing file is opened in read-only mode.
	\texttt{FileCreationAttr} is used to determine the MPI position part of the filename.
	If the file does not exist, an exception is thrown.

	\item \texttt{FAT\_READ\_MERGED}
	All existing files belonging to a single parallel run are opened simultaneously in read-only mode.
	Data from all files can be read transparently as if written to a single file.

\end{itemize}

\subsection{Closing Files}

After all file operations are finished and before opening or creating a new file,
already opened files must be closed by a call to \texttt{DataCollector::close}.
Otherwise, file information can be inconsistent and required data may not
be stored properly.

%-----------------------------------------------------------

\section{Datasets}

\emph{Datasets} are the general way for storing user data, e.g. simulation results or intermediate
systems states. They can be one-, two- or three-dimensional and each element can be
a basic type (e.g. \texttt{int}) or structured type (i.e. a \texttt{struct}).
Subclasses of the abstract class \texttt{CollectionType} are used to define the
\emph{Datatype} of a Dataset when storing data (or attributes to data).
The range of available types can be easily extended by defining a new subclass
(see \emph{include/basetypes} for a list of existing types).

Datasets are stored in the \emph{data} group of libSplash files and must be
related to a specific index beneath this group.

\subsection{Writing}
\label{lab:sdc:write}

To write data, use any of the \texttt{DataCollector::write} methods.
They require the used datatype, the number
of dimensions (\texttt{rank}, 1-3), the size of the actual data, buffers and offsets in
each dimension as a \texttt{Dimensions} object, the name for the dataset and a pointer
holding the data.
Any existing dataset in this group with the same is name is overwritten.

Please note that the user is responsible to pass a correct \texttt{CollectionType} to
any \texttt{write} call. Otherwise, user data may be interpreted incorrectly.

\subsection{Reading}

To read data, use any of the \texttt{DataCollector::read} methods.
They work similar to \texttt{DataCollector::write} but do not require
a \texttt{CollectionType} or \texttt{rank}, as these information are implicitly given from
the read dataset.
The destination buffer for reading must be allocated by the user. However, \texttt{read} methods
can be passed a \texttt{NULL} pointer to not read any data but only return the required dimensions of the
destination buffer.

\subsection{Appending}

Appending data is possible only for one-dimensional datasets.
It is achieved using any of the \texttt{DataCollector::append} methods.
In contrast to writing a dataset, any existing data remains unchanged and
new elements are appended at the end.
If the dataset for appending does not exist, it is created.

\subsection{Removing}

Datasets as well as whole program iterations can be removed from a file using\\ \texttt{DataCollector::remove}.
However, it may depend on the linked HDF5 library if file size actually decreases.

%-----------------------------------------------------------

\section{Attributes}

\emph{Attributes} are annotations to Datasets which can be used to store meta information.
Attributes can be of any \texttt{CollectionType} but must only contain a single element
of that type (in contrast to Datasets, which are multi-dimensional arrays).
At the moment, it is not possible to annotate Attributes at groups.

\subsection{Writing}

Attributes are written using the \texttt{DataCollector::writeAttribute} method.
It must be passed the location of the annotated Dataset (id and name) and the type and
name of the Attribute.

\subsection{Reading}

To read an Attribute, use the \texttt{DataCollector::readAttribute} method.
It must be passed the location of the annotated Dataset (id and name) and the
name of the Attribute. If no Attribute with this name and location exists, an
exception is thrown.

If the file has been opened for transparent merged read using \texttt{FAT\_READ\_MERGED},
additionally a MPI position can be defined to specify from which subfile to read
the required Attribute.
If this MPI position is set to \texttt{NULL}, the Attribute is read from the file with
MPI position (0, 0, 0).

\subsection{Global Attributes}

\emph{Global Attributes} are not specific to a single Dataset but belong to the whole
file or to each subfile, respectively.
To write and read Global Attributes, use\\ \texttt{DataCollector::writeGlobalAttribute} and\\
\texttt{DataCollector::readGlobalAttribute} methods.
Reading and writing is similar to normal Attributes, including the optional MPI position
when reading Global Attributes in \texttt{FAT\_READ\_MERGED} mode.

%-----------------------------------------------------------

\section{References}

References are links to another Dataset within one HDF5 file.
It can reference the whole Dataset as well as a user-defined subset, specified
by \emph{offset}, \emph{count} and \emph{stride}.
After a reference is created using \texttt{DataCollector::createReference},
it can be accessed like a normal Dataset.

%-----------------------------------------------------------

\chapter{DomainCollector}

The \texttt{DomainCollector} class extends \texttt{SerialDataCollector} with \emph{Domain}-specific operations.
A Domain is a logical data structure which can represent for example the simulation area of a iterative simulation program.
Each MPI process may write a sub-Domain (i.e. a sub-part) of the whole Domain area.
When accessing stored Domain data for analysis, it is possible to read data as stored in the HDF5 file as well as requesting
a part of the simulation Domain. This request is than handled transparently by libSplash.

Domain data can be of two types:
\begin{itemize}
	\item \textbf{GridType}
	This type is used for data stored as 1-3-dimensional arrays, such as fields or volumes where each element has a specified position within the domain grid.

	\item \textbf{PolyType}
	This type is used for unordered 1-dimensional data, e.g. particles within a volume.
\end{itemize}

\section{Writing Domains}

To write Domain data, use any of the \texttt{DomainCollector::writeDomain} calls.
Each requires information on the source data and buffer to read from, a name for the Dataset, the Datatype as well as
Domain information: the stored domain offset and area and the type of data stored (GridType or PolyType).

\section{Reading Domains}

To read Domain data, use \texttt{DomainCollector::readDomain}.
It must be passed the name and area of the requested Domain partition (sub-Domain).
The Domain type and a \texttt{DataContainer} holding the read data are returned.

When reading Domain data, read calls to multiple HDF5 Datasets may be necessary if the requested
sub-Domain spans multiple Datasets. Therefor, the returned DataContainer holds multiple
\texttt{DomainData} objects.
Each DomainData object stores a part of the requested data along with sub-Domain-specific information.
If GridType data is read, only a single DomainData object may be returned as data from multiple Datasets can
be transparently combined into one new array.
Otherwise, when reading PolyType data, the DataContainer is likely to hold multiple DomainData objects, one for
each physical read.
These DomainData object can be queried successively or by their 1-3.dimensional index (\texttt{DataContainer::get} and \texttt{DataContainer::getIndex}).
Additionally, all elements from all DomainData objects within one DataContainer can be queried continuously using \texttt{DataContainer::getNumElements} and
\texttt{DataContainer::getElement}.

\section{Appending Domains}

Appending Domain data follows the same restrictions as appending normal Datasets (see \ref{lab:sdc:write}).

%-----------------------------------------------------------

\chapter{Misc}

\section{splashtools}

\textbf{splashtools} is a convenience tool to check/modify/... HDF5 files created with libSplash.
Features include:
\begin{itemize}
	\item Transparently delete timesteps in all HDF5 files belonging to a single run.
	\item Check files for syntactic and semantic consistency.
\end{itemize}
Run \command{splashtools --help} for a complete list of all current features.


\section{Tests}

The libSplash repository contains tests for self-testing the library. They can be found in the
\command{tests} subdirectory. To build the tests, move to the \command{tests} subdirectory and execute\\
\command{mkdir build; cd build; cmake ..; make}.\\
From the \command{tests} directory, all tests can be run using the \command{run\_tests.sh} shell script.
To build the tests, \textbf{cppunit} and \textbf{OpenMPI} must be installed.

%-----------------------------------------------------------



\end{document}
